\section{Requisitos}

\subsection{Requisitos no funcionales}

Las características de calidad corresponden a ISO/IEC 25010:2023.

\begin{description}
    \item[RNF-01. Capacidad documental.] Dimensionar el sistema para 10 a
    12 estudiantes por periodo y hasta tres PDF por estudiante, con un límite
    de 20 MB por archivo. El máximo previsto de archivos cargados es de
    720 MB por periodo (12 estudiantes por 3 archivos por 20 MB), a lo que
    se suma el espacio de las actas generadas. La nómina del segundo periodo
    2024--2025 tuvo 12 estudiantes inscritos y 10 presentaron. Verificar la
    capacidad con ese volumen y probar el límite de tamaño por archivo.
    \textbf{Característica.} Eficiencia de desempeño (capacidad).

    \item[RNF-02. Disponibilidad.] Alcanzar una disponibilidad mensual de al
    menos 99,5\% fuera de la semana de predefensas. Durante esa semana,
    recuperar las funciones críticas en un máximo de 60 minutos consecutivos
    por interrupción. Verificar ambos criterios con registros de monitoreo.
    Si el sistema no está disponible, el tribunal puede continuar la
    predefensa y registrar las calificaciones después.
    \textbf{Característica.} Fiabilidad (disponibilidad y recuperabilidad).

    \item[RNF-03. Acceso desde computadora y celular.] Permitir la carga de
    documentos, la revisión de informes y la consulta del estado desde
    ambos tipos de dispositivo, según el rol. Verificar esas operaciones
    en 2 entornos: un navegador de escritorio y uno móvil.
    \textbf{Característica.} Flexibilidad (adaptabilidad).

    \item[RNF-04. Conservación documental.] Conservar permanentemente las
    actas y las notas del expediente académico. Aplicar a los demás
    documentos el plazo de la tabla institucional de conservación
    documental, pendiente de confirmar. El proceso automático de limpieza
    debe eliminar 0 actas y 0 registros de notas. Verificarlo ejecutando
    ese proceso sobre expedientes de prueba.
    \textbf{Característica.} Seguridad (integridad).

    \item[RNF-05. Concurrencia.] Soportar 30 usuarios simultáneos en una
    prueba de carga que combine entrega de archivos, revisión de informes y
    consulta del estado, con 0 errores atribuibles a saturación durante la
    prueba. Los 30 usuarios son una meta preliminar de diseño basada en la
    estimación de pocas decenas de usuarios durante las predefensas, que se realizan en la semana 19 del décimo semestre, y en los días previos de entrega.
    \textbf{Característica.} Eficiencia de desempeño (capacidad).

    \item[RNF-06. Confidencialidad.] Permitir 0 accesos no autorizados a
    informes y calificaciones. Verificar, con al menos 2 expedientes y
    2 tribunales de prueba, que un estudiante no pueda acceder a un
    expediente ajeno y que un tribunal no pueda acceder a expedientes
    asignados al otro tribunal.
    \textbf{Característica.} Seguridad (confidencialidad).
\end{description}
